% تمرین نمونه است تا مخزن از روز نخست خالی نباشد؛ پیش از استفاده بازبینی‌اش کنید.
\documentclass{../is-assignment}

\عنوان{\متن‌لاتین{Crypto Toolbox}}
\ترم{پاییز ۱۴۰۵}

\begin{document}

\عنوان‌ساز

\فهرست‌مطالب

\قسمت{مقدمه}

در این تمرین ابزارهای رمزنگاری‌ای را که در کلاس دیدیم روی یک پرونده‌ی واقعی به کار می‌بریم:
رمزگذاری متقارن، تشخیص دستکاری، و امضای دیجیتال.
هدف این است که تفاوت «محرمانه ماندن» با «دستکاری نشدن» و با «قابل اثبات بودن فرستنده» را با دست خودتان ببینید،
نه اینکه تعریفشان را بنویسید.

استفاده از \متن‌لاتین{openssl} یا کتابخانه‌ی رمزنگاری هر زبانی که می‌پسندید آزاد است.
پیاده‌سازی \متن‌سیاه{خودِ} الگوریتم‌ها خواسته نشده و امتیازی هم ندارد.

\قسمت{محرمانگی}

پرونده‌ای با دست‌کم ۱۰ کیلوبایت متن بسازید و آن را با \متن‌لاتین{AES-256-GCM} رمز کنید.

\شروع{شمارش}
\فقره فرمان یا کدی که استفاده کرده‌اید را بیاورید.
\فقره کلید و \متن‌لاتین{nonce} را چگونه ساختید؟ چرا \متن‌لاتین{nonce} نباید با همان کلید دوباره استفاده شود؟
\فقره یک بایت از متن رمزشده را عوض کنید و دوباره رمزگشایی کنید. چه اتفاقی می‌افتد و چرا؟
\پایان{شمارش}

\نمره{۴}

\قسمت{یکپارچگی}

\شروع{شمارش}
\فقره برای همان پرونده \متن‌لاتین{SHA-256} و \متن‌لاتین{HMAC-SHA-256} را حساب کنید.
\فقره فرض کنید پرونده و مقدار حساب‌شده هر دو از یک مسیر ناامن می‌گذرند.
      کدام یک جلوی دستکاری را می‌گیرد و کدام نه؟ با یک مثال \متن‌سیاه{عملی} نشان دهید، یعنی خودتان نقش مهاجم را بازی کنید.
\پایان{شمارش}

\نمره{۳}

\قسمت{اصالت}

\شروع{شمارش}
\فقره یک جفت کلید بسازید، پرونده را امضا کنید و امضا را با کلید عمومی وارسی کنید.
\فقره چرا \متن‌لاتین{HMAC} برای اثبات نزد شخص سوم کافی نیست ولی امضا هست؟
\پایان{شمارش}

\نمره{۳}

\شروع{امتیازی}
گواهی‌نامه‌ی یک وب‌گاه را با \متن‌لاتین{openssl s\_client} بگیرید و زنجیره‌ی اعتماد آن را تا ریشه دنبال کنید.
بگویید مرورگر شما دقیقا به چه چیزی اعتماد می‌کند و اگر یک حلقه‌ی میانی از دست برود چه می‌شود.
\پایان{امتیازی}

\قسمت{تحویل}

گزارشی کوتاه (حداکثر چهار صفحه) به همراه فرمان‌ها، خروجی‌ها و پاسخ پرسش‌ها.
خروجی‌ها باید از اجرای خودتان باشد؛ خروجی کپی‌شده بدون اجرا، تحویل به شمار نمی‌آید.

\قوانین

\پایان‌ساز

\end{document}
